\begin{proof}
\textbf{Step 1. A non‑negative quadratic.}  
For a real parameter \(t\) define  
\[
f(t)=\sum_{i=1}^{n}\bigl(a_{i}-t\,b_{i}\bigr)^{2}.
\tag{1}
\]
Each summand is a square, hence \(f(t)\ge 0\) for all \(t\in\mathbb{R}\).

\textbf{Step 2. Expanding \(f(t)\).}  
Using \((x-y)^{2}=x^{2}-2xy+y^{2}\) we obtain  
\[
(a_{i}-t b_{i})^{2}=a_{i}^{2}-2t\,a_{i}b_{i}+t^{2}b_{i}^{2}.
\]
Summing over \(i\) gives  
\[
\begin{aligned}
f(t)&=\sum_{i=1}^{n}a_{i}^{2}
      -2t\sum_{i=1}^{n}a_{i}b_{i}
      +t^{2}\sum_{i=1}^{n}b_{i}^{2}   \\
    &=A-2Ct+Bt^{2},
\end{aligned}
\tag{2}
\]
where  
\[
A:=\sum_{i=1}^{n}a_{i}^{2}\ge 0,\qquad
B:=\sum_{i=1}^{n}b_{i}^{2}\ge 0,\qquad
C:=\sum_{i=1}^{n}a_{i}b_{i}.
\]

\textbf{Step 3. Discriminant condition.}  
A quadratic polynomial \(qt^{2}+pt+r\) satisfies \(qt^{2}+pt+r\ge0\) for every real \(t\) iff either  

* \(q>0\) and its discriminant is non‑positive: \(p^{2}-4qr\le0\); or  
* \(q=0\) and the resulting linear (or constant) function is non‑negative for all \(t\).

In \eqref{2} the leading coefficient is \(q=B\).

\begin{itemize}
\item If \(B>0\), then \(p=-2C\) and \(r=A\). The discriminant condition yields  
\[
(-2C)^{2}-4BA\le0\;\Longrightarrow\;C^{2}\le AB,
\]
i.e.  
\[
\Bigl(\sum_{i=1}^{n}a_{i}b_{i}\Bigr)^{2}
\le
\Bigl(\sum_{i=1}^{n}a_{i}^{2}\Bigr)
\Bigl(\sum_{i=1}^{n}b_{i}^{2}\Bigr).
\tag{3}
\]

\item If \(B=0\), then every \(b_{i}=0\). Hence \(C=0\) and \eqref{3} reads \(0\le0\), which is true. The same argument applies when \(A=0\).
\end{itemize}
Thus \eqref{3} holds in all cases.

\textbf{Step 4. Equality case.}  
Equality in \eqref{3} occurs precisely when the discriminant is zero (or when \(B=0\) or \(A=0\)).

\begin{itemize}
\item \emph{Generic case \(B>0\).}  
If the discriminant is zero, then \(C^{2}=AB\) and the quadratic can be written as  
\[
f(t)=B\bigl(t-\lambda\bigr)^{2},\qquad \lambda:=\frac{C}{B}.
\]
Hence
\[
\sum_{i=1}^{n}\bigl(a_{i}-\lambda b_{i}\bigr)^{2}=0,
\]
which forces each term to vanish:
\[
a_{i}=\lambda b_{i}\qquad\text{for all }i.
\]

\item \emph{Degenerate case \(B=0\).}  
Then \(b_{i}=0\) for all \(i\). The inequality becomes \(0\le0\) and equality holds for every choice of \((a_{i})\). Since the zero vector is a scalar multiple of any vector, the condition “there exists \(\lambda\) with \(a_{i}=\lambda b_{i}\) for all \(i\)” is still satisfied (any \(\lambda\) works). An analogous argument applies when \(A=0\).
\end{itemize}
Consequently, equality holds \emph{iff} the two vectors are linearly dependent, i.e. there exists \(\lambda\in\mathbb{R}\) such that \(a_{i}=\lambda b_{i}\) for every \(i\) (including the degenerate case where one of the vectors is the zero vector). \(\square\)
\end{proof}